\subsection{Наблюдавшиеся компоненты RAM на общей нагрузке}
\label{sec:ram500}
Отдельная диагностическая серия RAM500 выполнена 16 сентября 2026 года на тех же экземплярах Mega 2560 и ESP32-C3. Для каждой из пяти зафиксированных моделей сделаны три свежие загрузки прошивки; каждая выполняет три прохода по общим 500 подготовленным входам. На каждой плате получены 22\,500 индивидуальных сравнений с сохранёнными эталонными предсказаниями, без расхождений. Это повторное исполнение 500 входов, а не 22\,500 различных потоков и не отсутствие ошибок относительно истинных меток. Модели и входы хранятся во Flash, очередной вход копируется в рабочий буфер. USB-тестер в измерении RAM не участвует.

После компоновки статические данные и BSS диагностической программы занимают 335 байт SRAM на Mega и 15\,312 байт DRAM на C3 для каждой модели. Эти размеры включают измерительную программу и не заменяют размеры исторических или энергетических приложений. На Mega watermark, то есть граница затронутой области памяти, регистрирует наблюдавшуюся глубину стека в области загрузки входа, предсказания и сверки с эталоном, включая возникающие прерывания. На C3 используется lifetime watermark задачи \texttt{loop}: он отражает её предшествовавшую историю с момента создания, включая подготовку и обмен служебными сообщениями~\cite{EspressifFreeRTOS447}. Поэтому столбцы табл.~\ref{tab:ram500} имеют разные области измерения.

\begin{table}[!htbp]
\centering\small\setlength{\tabcolsep}{8pt}
\caption{Наблюдавшееся использование стека диагностической программы, байт. Mega: одинаковое значение во всех трёх загрузках. C3: занятость стека задачи по lifetime watermark отдельно для трёх загрузок; во всех запусках максимум уже достигнут до рабочих проходов. Значения C3 не определяют стек только инференса.}
\label{tab:ram500}
\begin{tabular}{lrrrr}
\toprule
& Mega & \multicolumn{3}{c}{ESP32-C3: загрузка}\\
\cmidrule(lr){3-5}
Модель & стек & 1 & 2 & 3\\
\midrule
KAN, коэффициенты & 93 & 1584 & 1584 & 1584\\
KAN, LUT & 80 & 1504 & 1584 & 1584\\
MLP16 & 95 & 1584 & 1584 & 1584\\
Многослойная KAN & 208 & 1504 & 1504 & 1504\\
DT5 & 46 & 1584 & 1584 & 1504\\
\bottomrule
\end{tabular}
\end{table}

На C3 зарезервированы 8192 байта стека задачи \texttt{loop}; незатронутый остаток составляет не менее 6608 байт. Watermark до и после нагрузки совпадает во всех пятнадцати запусках: инференс не превысил ранее достигнутую границу, но это не означает отсутствия использования стека. Различие 80 байт между отдельными загрузками нельзя приписывать модели. Снимки свободной внутренней 8-битно адресуемой heap~\cite{EspressifHeap447} до и после равны 309\,800 байт; наибольший свободный блок --- 286\,708 байт. Сумма lifetime-минимумов регионов равна 304\,820 байт и не является измеренным одновременным минимумом всей heap. Неизменность выборочных показаний не исключает кратковременных выделений между ними. Резерв стека уже входит в выделения heap; IRAM, статические данные, резервированные и фактически затронутые области нельзя суммировать как независимые слагаемые единого Peak RAM.

Положительные контроли выполнены после сохранения основных показаний и проходят во всех запусках. На Mega контрольное выделение изменяет границу heap на 130 байт с последующим возвратом; в рабочих проходах счётчики перехваченных \texttt{malloc}, \texttt{calloc}, \texttt{realloc} и \texttt{free} равны нулю. На C3 запрос 512 байт уменьшает свободную heap на 528 байт, затем значение восстанавливается; watermark отдельной контрольной задачи изменяется на 676 байт. Эти проверки подтверждают чувствительность каналов, а не калибруют глобальный максимум памяти.

Серия устанавливает успешное исполнение заданной нагрузки при наблюдавшемся запасе памяти. Она не даёт общего worst-case Peak RAM для всех входов, задач RTOS и режимов приложения. Отдельный эксперимент со свежей изолированной задачей потребовался бы для сравнения стека только инференса на C3; такой результат не входит в область заявленных измерений.
